\documentclass[10pt,a4paper]{amsart}
\usepackage[margin=19mm]{geometry}
\usepackage{amsmath,amssymb,mathtools}
\usepackage[scr=rsfs,cal=euler]{mathalfa}
\usepackage{xfrac}
\usepackage[unicode,hidelinks,bookmarks=false]{hyperref}
\newcommand{\ie}{i.e.\ }
\newcommand{\cf}{cf.\ }
\newcommand{\resp}{respectively\ }
\newcommand{\Subsection}{\subsection}
\providecommand{\nobreakdash}{\nobreak\hbox{-}}
\setlength{\emergencystretch}{2em}
\multlinegap0pt
\setcounter{MaxMatrixCols}{30}%
\providecommand{\U}[1]{\protect\rule{.1in}{.1in}}
\newtheorem{theorem}{Theorem}
\newtheorem{acknowledgement}{Acknowledgement}
\newtheorem{algorithm}{Algorithm}
\newtheorem{axiom}{Axiom}
\newtheorem{case}{Case}
\newtheorem{claim}{Claim}
\newtheorem{conclusion}{Conclusion}
\newtheorem{condition}{Condition}
\newtheorem{conjecture}{Conjecture}
\newtheorem{corollary}{Corollary}
\newtheorem{criterion}{Criterion}
\newtheorem{definition}{Definition}
\newtheorem{example}{Example}
\newtheorem{exercise}{Exercise}
\newtheorem{lemma}{Lemma}
\newtheorem{notation}{Notation}
\newtheorem{problem}{Problem}
\newtheorem{proposition}{Proposition}
\newtheorem{remark}{Remark}
\newtheorem{solution}{Solution}
\newtheorem{summary}{Summary}
\numberwithin{equation}{section}
\begin{document}
\title[The Fractional Half Dirac Operator Over Sobolev Spaces]{The Fractional Half Dirac Operator Over Sobolev Spaces}
\author{Dejenie A. Lakew}
\date{}
\maketitle
\noindent\emph{Source and licence.} The Fractional Half Dirac Operator Over Sobolev Spaces. Version arXiv:2608.02871v1 (CC BY 4.0). This material is distributed under the Creative Commons Attribution 4.0 International licence, http://creativecommons.org/licenses/by/4.0/, as recorded in the arXiv OAI licence metadata for this exact version and shown on the abstract page https://arxiv.org/abs/2608.02871v1. Full rights notice: licenses/arxiv-2608.02871-CC-BY-4.0.txt.\par\smallskip
\noindent\textit{Editorial scope.} Counted unit: the original \texttt{theorem} environment \texttt{-} at source lines 1826--1887 together with its complete proof at lines 1889--2339; the delivered document keeps source lines 1819--2340 so that every referenced label and every citation resolves inside it. Native editable source is the paper's own concatenated TeX; the AMS wrapper is editorial formatting only. No publisher class, upstream script or unrelated artwork is redistributed.
\section{Boundary Value Problems and Decompositions}

In this section we consider first and second order boundary value problems and
see how the solutions are simply the orthogonal sums of parts that evolve from
boundary values of the solution and interior values of the derivative of the
solution to the given order. We start with the first order Cauchy problem.


\begin{theorem}
\begin{proposition}
Let $f\in%
%TCIMACRO{\tciLaplace}%
%BeginExpansion
\mathcal{L}%
%EndExpansion
^{2}\left(  \Omega\right)  $,$g\in%
%TCIMACRO{\tciLaplace}%
%BeginExpansion
\mathcal{L}%
%EndExpansion
^{2}\left(  \partial\Omega\right)  $ The first Caputo fractional $1/2$ order
Cauchy problem%
\begin{equation}
\left\{
\begin{array}
[c]{cc}%
D^{\frac{1}{2}}u=f & \text{in }\Omega\\
u=g & \text{on }\partial\Omega
\end{array}
\right.
\end{equation}
has a solution $u\in W^{\frac{1}{2},2}\left(  \Omega\right)  $ such that
\[
u=\left[  u\right]  _{f}\uplus\left[  u\right]  _{g}%
\]
where $\left[  u\right]  _{f}$ is the part of the solution that evolves from
$f $ and $\left[  u\right]  _{g}$ is the part of the solution that evolves
from $g $ with the following properties

$(i)$
\[
\left\langle \left[  u\right]  _{f},\text{ }\left[  u\right]  _{g}%
\right\rangle _{W^{\frac{1}{2},2}\left(  \Omega\right)  }=0
\]
$\left(  ii\right)  $%
\[
\parallel u\parallel_{W^{\frac{1}{2},2}\left(  \Omega\right)  }^{2}%
=\parallel\left[  u\right]  _{f}\parallel_{W^{\frac{1}{2},2}\left(
\Omega\right)  }^{2}+\parallel\left[  u\right]  _{g}\parallel_{W^{\frac{1}%
{2},2}\left(  \Omega\right)  }^{2}%
\]

\end{proposition}

$\left(  iii\right)  $
\[
\left\langle u,\text{ }\left[  u\right]  _{f}\right\rangle _{W^{\frac{1}{2}%
,2}\left(  \Omega\right)  }=\parallel\left[  u\right]  _{f}\parallel
_{W^{\frac{1}{2},2}\left(  \Omega\right)  }^{2}%
\]


$\left(  iv\right)  $
\[
\left\langle u,\text{ }\left[  u\right]  _{g}\right\rangle _{W^{\frac{1}{2}%
,2}\left(  \Omega\right)  }=\parallel\left[  u\right]  _{g}\parallel
_{W^{\frac{1}{2},2}\left(  \Omega\right)  }^{2}%
\]

\end{theorem}

\begin{proof}
For two functions $f,$ $g\in C^{1}\left(  \Omega\right)  $, integration by
parts provide%

\[%
%TCIMACRO{\dint \limits_{\Omega}}%
%BeginExpansion
{\displaystyle\int\limits_{\Omega}}
%EndExpansion
f\left(  y-x\right)  D_{y}^{\frac{1}{2}}g(y)d\Omega_{y}=%
%TCIMACRO{\dint \limits_{\partial\Omega}}%
%BeginExpansion
{\displaystyle\int\limits_{\partial\Omega}}
%EndExpansion
f\left(  y-x\right)  \upsilon_{\frac{1}{2}}\left(  y\right)  g\left(
y\right)  d\partial\Omega_{y}-%
%TCIMACRO{\dint \limits_{\Omega}}%
%BeginExpansion
{\displaystyle\int\limits_{\Omega}}
%EndExpansion
D_{y}^{\frac{1}{2}}f\left(  y-x\right)  g\left(  y\right)  d\Omega_{y}.
\]
Then taking $f=\Gamma_{\frac{1}{2}}$ and $g=u$, we have
\[%
%TCIMACRO{\dint \limits_{\Omega}}%
%BeginExpansion
{\displaystyle\int\limits_{\Omega}}
%EndExpansion
\Gamma_{\frac{1}{2}}\left(  y-x\right)  D_{y}^{\frac{1}{2}}u(y)d\Omega_{y}=%
%TCIMACRO{\dint \limits_{\partial\Omega}}%
%BeginExpansion
{\displaystyle\int\limits_{\partial\Omega}}
%EndExpansion
\Gamma_{\frac{1}{2}}\left(  y-x\right)  \upsilon_{\frac{1}{2}}\left(
y\right)  u\left(  y\right)  d\partial\Omega_{y}-%
%TCIMACRO{\dint \limits_{\Omega}}%
%BeginExpansion
{\displaystyle\int\limits_{\Omega}}
%EndExpansion
D_{y}^{\frac{1}{2}}\Gamma_{\frac{1}{2}}\left(  y-x\right)  u\left(  y\right)
d\Omega_{y}%
\]
where $\Gamma_{\frac{1}{2}}$ is the fundamental solution to the fractional
half Dirac operator $D^{\frac{1}{2}}$.

But
\[
D_{y}^{\frac{1}{2}}\Gamma_{\frac{1}{2}}\left(  y-x\right)  =\delta\left(
y-x\right)
\]
the Kronecker delta function and thus, \ %

\begin{align*}%
%TCIMACRO{\dint \limits_{\Omega}}%
%BeginExpansion
{\displaystyle\int\limits_{\Omega}}
%EndExpansion
\Gamma_{\frac{1}{2}}\left(  y-x\right)  D^{\frac{1}{2}}u(y)d\Omega_{y}  & =%
%TCIMACRO{\dint \limits_{\partial\Omega}}%
%BeginExpansion
{\displaystyle\int\limits_{\partial\Omega}}
%EndExpansion
\Gamma_{\frac{1}{2}}\left(  y-x\right)  \upsilon_{\frac{1}{2}}\left(
y\right)  u\left(  y\right)  d\partial\Omega_{y}-%
%TCIMACRO{\dint \limits_{\Omega}}%
%BeginExpansion
{\displaystyle\int\limits_{\Omega}}
%EndExpansion
\delta\left(  y-x\right)  u\left(  y\right)  d\Omega_{y}\\
& =%
%TCIMACRO{\dint \limits_{\partial\Omega}}%
%BeginExpansion
{\displaystyle\int\limits_{\partial\Omega}}
%EndExpansion
\Gamma_{\frac{1}{2}}\left(  y-x\right)  \upsilon_{\frac{1}{2}}\left(
y\right)  u\left(  y\right)  d\partial\Omega_{y}-u(x)
\end{align*}


i.e.,
\begin{equation}
u\left(  x\right)  =%
%TCIMACRO{\dint \limits_{\partial\Omega}}%
%BeginExpansion
{\displaystyle\int\limits_{\partial\Omega}}
%EndExpansion
\Gamma_{\frac{1}{2}}\left(  y-x\right)  \upsilon_{\frac{1}{2}}\left(
y\right)  u\left(  y\right)  d\partial\Omega_{y}-%
%TCIMACRO{\dint \limits_{\Omega}}%
%BeginExpansion
{\displaystyle\int\limits_{\Omega}}
%EndExpansion
\Gamma_{\frac{1}{2}}\left(  y-x\right)  D^{\frac{1}{2}}u(y)d\Omega
_{y}\label{borel-pompeiu}%
\end{equation}
which provides the integral representation of the solution $u$ to the BVP
given by
\[
u\left(  x\right)  =%
%TCIMACRO{\dint \limits_{\partial\Omega}}%
%BeginExpansion
{\displaystyle\int\limits_{\partial\Omega}}
%EndExpansion
\Gamma_{\frac{1}{2}}\left(  y-x\right)  \upsilon_{\frac{1}{2}}\left(
y\right)  g(y)d\partial\Omega_{y}+\left(  -%
%TCIMACRO{\dint \limits_{\Omega}}%
%BeginExpansion
{\displaystyle\int\limits_{\Omega}}
%EndExpansion
\Gamma_{\frac{1}{2}}\left(  y-x\right)  f(y)d\Omega_{y}\right)  .
\]
We will see that this sum is in fact an orthogonal sum $\uplus$. \ 

Because the solution
\[
u\in W^{\frac{1}{2},2}\left(  \Omega\right)  =A^{\frac{1}{2},2}\left(
\Omega\right)  \oplus D^{\frac{1}{2}}\left(  W_{0}^{\frac{1}{2},2}\left(
\Omega\right)  \right)
\]
has a unique decomposition as sum of components from the two sub spaces,
$A^{\frac{1}{2},2}\left(  \Omega\right)  $ and $D^{\frac{1}{2}}\left(
W_{0}^{\frac{3}{2},2}\left(  \Omega\right)  \right)  $ and the first integral
\[%
%TCIMACRO{\dint \limits_{\partial\Omega}}%
%BeginExpansion
{\displaystyle\int\limits_{\partial\Omega}}
%EndExpansion
\Gamma_{\frac{1}{2}}\left(  y-x\right)  \upsilon_{\frac{1}{2}}\left(
y\right)  g(y)d\partial\Omega_{y}%
\]
is monogenic over $\Omega$ and hence an element of
\[
KerD^{\frac{1}{2}}\cap%
%TCIMACRO{\tciLaplace}%
%BeginExpansion
\mathcal{L}%
%EndExpansion
^{2}\left(  \Omega\right)  =A^{\frac{1}{2},2}\left(  \Omega\right)  .
\]
We need to verify that the second integral%

\[
\left(  -%
%TCIMACRO{\dint \limits_{\Omega}}%
%BeginExpansion
{\displaystyle\int\limits_{\Omega}}
%EndExpansion
\Gamma_{\frac{1}{2}}\left(  y-x\right)  f(y)d\Omega_{y}\right)  \in
D^{\frac{1}{2}}\left(  W_{0}^{\frac{1}{2},2}\left(  \Omega\right)  \right)
\]
as well. That is, there is a function $\xi\in W_{0}^{\frac{3}{2},2}\left(
\Omega\right)  $ so that%

\[
-%
%TCIMACRO{\dint \limits_{\Omega}}%
%BeginExpansion
{\displaystyle\int\limits_{\Omega}}
%EndExpansion
\Gamma_{\frac{1}{2}}\left(  y-x\right)  f(y)d\Omega_{y}=D^{\frac{1}{2}}%
\xi\left(  x\right)  \text{ \ \ with \ \ }\xi_{\mid\partial\Omega}=0.
\]
Clearly from the fact that \ \ %

\[
u\left(  x\right)  =%
%TCIMACRO{\dint \limits_{\partial\Omega}}%
%BeginExpansion
{\displaystyle\int\limits_{\partial\Omega}}
%EndExpansion
\Gamma_{\frac{1}{2}}\left(  y-x\right)  \upsilon_{\frac{1}{2}}\left(
y\right)  g(y)d\partial\Omega_{y}+\left(  -%
%TCIMACRO{\dint \limits_{\Omega}}%
%BeginExpansion
{\displaystyle\int\limits_{\Omega}}
%EndExpansion
\Gamma_{\frac{1}{2}}\left(  y-x\right)  f(y)d\Omega_{y}\right)  \text{ \ and
\ }u_{\mid\partial\Omega}=g
\]
we have, the trace $\tau$ of $u$ on $\partial\Omega$%

\begin{align*}
\tau u_{\mid\partial\Omega}  & =\tau\left(
%TCIMACRO{\dint \limits_{\partial\Omega}}%
%BeginExpansion
{\displaystyle\int\limits_{\partial\Omega}}
%EndExpansion
\Gamma_{\frac{1}{2}}\left(  y-x\right)  \upsilon_{\frac{1}{2}}\left(
y\right)  g(y)d\partial\Omega_{y}+\left(  -%
%TCIMACRO{\dint \limits_{\Omega}}%
%BeginExpansion
{\displaystyle\int\limits_{\Omega}}
%EndExpansion
\Gamma_{\frac{1}{2}}\left(  y-x\right)  f(y)d\Omega_{y}\right)  \right)
_{\mid\partial\Omega}\\
& =\tau\left(
%TCIMACRO{\dint \limits_{\partial\Omega}}%
%BeginExpansion
{\displaystyle\int\limits_{\partial\Omega}}
%EndExpansion
\Gamma_{\frac{1}{2}}\left(  y-x\right)  \upsilon_{\frac{1}{2}}\left(
y\right)  g(y)d\partial\Omega_{y}\right)  _{\mid\partial\Omega}+\tau\left(  -%
%TCIMACRO{\dint \limits_{\Omega}}%
%BeginExpansion
{\displaystyle\int\limits_{\Omega}}
%EndExpansion
\Gamma_{\frac{1}{2}}\left(  y-x\right)  f(y)d\Omega_{y}\right)  _{\mid
\partial\Omega}\\
& =g
\end{align*}
%

\[
\Longrightarrow\text{ \ \ \ \ }\tau\left(  -%
%TCIMACRO{\dint \limits_{\Omega}}%
%BeginExpansion
{\displaystyle\int\limits_{\Omega}}
%EndExpansion
\Gamma_{\frac{1}{2}}\left(  y-x\right)  f(y)d\Omega_{y}\right)  =\left(  -%
%TCIMACRO{\dint \limits_{\Omega}}%
%BeginExpansion
{\displaystyle\int\limits_{\Omega}}
%EndExpansion
\Gamma_{\frac{1}{2}}\left(  y-x\right)  f(y)d\Omega_{y}\right)  _{\mid
\partial\Omega}=0\text{.}%
\]


\ \ 

Thus $\xi_{\mid\partial\Omega}=0$. In an analogous manner of the integral
representation of the solution $u$, we get the integral representation of
$\xi$ as

\
\begin{align*}
\xi\left(  x\right)   & =-%
%TCIMACRO{\dint \limits_{\Omega}}%
%BeginExpansion
{\displaystyle\int\limits_{\Omega}}
%EndExpansion
\Gamma_{\frac{1}{2}}\left(  z-x\right)
%TCIMACRO{\dint \limits_{\Omega}}%
%BeginExpansion
{\displaystyle\int\limits_{\Omega}}
%EndExpansion
\Gamma_{\frac{1}{2}}\left(  y-z\right)  f(y)d\Omega_{y}d\Omega_{z}\\
& =-%
%TCIMACRO{\dint \limits_{\Omega}}%
%BeginExpansion
{\displaystyle\int\limits_{\Omega}}
%EndExpansion%
%TCIMACRO{\dint \limits_{\Omega}}%
%BeginExpansion
{\displaystyle\int\limits_{\Omega}}
%EndExpansion
\Gamma_{\frac{1}{2}}\left(  z-x\right)  \Gamma_{\frac{1}{2}}\left(
y-z\right)  f(y)d\Omega_{y}d\Omega_{z}\in W_{0}^{1,2}\left(  \Omega\right)
\text{.}%
\end{align*}
Therefore, there is a \ %

\[
\xi\left(  x\right)  =-%
%TCIMACRO{\dint \limits_{\Omega}}%
%BeginExpansion
{\displaystyle\int\limits_{\Omega}}
%EndExpansion%
%TCIMACRO{\dint \limits_{\Omega}}%
%BeginExpansion
{\displaystyle\int\limits_{\Omega}}
%EndExpansion
\Gamma_{\frac{1}{2}}\left(  z-x\right)  \Gamma_{\frac{1}{2}}\left(
y-z\right)  f(y)d\Omega_{y}d\Omega_{z}\in W_{0}^{1,2}\left(  \Omega\right)
\]
so that \ \ \ %

\begin{align*}
D_{x}^{\frac{1}{2}}\xi\left(  x\right)   & =D_{x}^{\frac{1}{2}}\left(  -%
%TCIMACRO{\dint \limits_{\Omega}}%
%BeginExpansion
{\displaystyle\int\limits_{\Omega}}
%EndExpansion%
%TCIMACRO{\dint \limits_{\Omega}}%
%BeginExpansion
{\displaystyle\int\limits_{\Omega}}
%EndExpansion
\Gamma_{\frac{1}{2}}\left(  z-x\right)  \Gamma_{\frac{1}{2}}\left(
y-z\right)  f(y)d\Omega_{y}d\Omega_{z}\right) \\
& =-%
%TCIMACRO{\dint \limits_{\Omega}}%
%BeginExpansion
{\displaystyle\int\limits_{\Omega}}
%EndExpansion
\Gamma_{\frac{1}{2}}\left(  y-x\right)  f(y)d\Omega_{y}\text{.}%
\end{align*}
Hence \ %

\begin{align*}
u(x)  & =%
%TCIMACRO{\dint \limits_{\partial\Omega}}%
%BeginExpansion
{\displaystyle\int\limits_{\partial\Omega}}
%EndExpansion
\Gamma_{\frac{1}{2}}\left(  y-x\right)  \upsilon_{\frac{1}{2}}\left(
y\right)  g(y)d\partial\Omega_{y}\uplus\left(  -%
%TCIMACRO{\dint \limits_{\Omega}}%
%BeginExpansion
{\displaystyle\int\limits_{\Omega}}
%EndExpansion
\Gamma_{\frac{1}{2}}\left(  y-x\right)  f(y)d\Omega_{y}\right) \\
& =\left[  u\right]  _{g}\uplus\left[  u\right]  _{f}%
\end{align*}
with%

\[
\left[  u\right]  _{g}=%
%TCIMACRO{\dint \limits_{\partial\Omega}}%
%BeginExpansion
{\displaystyle\int\limits_{\partial\Omega}}
%EndExpansion
\Gamma_{\frac{1}{2}}\left(  y-x\right)  \upsilon_{\frac{1}{2}}\left(
y\right)  g(y)d\partial\Omega_{y}\text{ \ \ and \ \ }\left[  u\right]  _{f}=-%
%TCIMACRO{\dint \limits_{\Omega}}%
%BeginExpansion
{\displaystyle\int\limits_{\Omega}}
%EndExpansion
\Gamma_{\frac{1}{2}}\left(  y-x\right)  f(y)d\Omega_{y}\text{.}%
\]


\ 

$\left(  i\right)  $ Then since $W^{\frac{1}{2},2}\left(  \Omega\right)  $ is
an inner product space of an orthogonal decomposition we have%

\[
\left\langle \left[  u\right]  _{g},\left[  u\right]  _{f}\right\rangle
=\left\langle
%TCIMACRO{\dint \limits_{\partial\Omega}}%
%BeginExpansion
{\displaystyle\int\limits_{\partial\Omega}}
%EndExpansion
\Gamma_{\frac{1}{2}}\left(  y-x\right)  \upsilon_{\frac{1}{2}}\left(
y\right)  g(y)d\partial\Omega_{y},-%
%TCIMACRO{\dint \limits_{\Omega}}%
%BeginExpansion
{\displaystyle\int\limits_{\Omega}}
%EndExpansion
\Gamma_{\frac{1}{2}}\left(  y-x\right)  f(y)d\Omega_{y}\right\rangle
_{W^{\frac{1}{2},2}\left(  \Omega\right)  }=0\text{.}%
\]


\ 

$\left(  ii\right)  $ \ From the fact that the sum is an orthogonal sum, the
components obey the parallelogram law

\
\begin{align*}
& \parallel u\parallel_{W^{\frac{1}{2},2}\left(  \Omega\right)  }%
^{2}=\parallel%
%TCIMACRO{\dint \limits_{\partial\Omega}}%
%BeginExpansion
{\displaystyle\int\limits_{\partial\Omega}}
%EndExpansion
\Gamma_{\frac{1}{2}}\left(  y-x\right)  \upsilon_{\frac{1}{2}}\left(
y\right)  g(y)d\partial\Omega_{y}\parallel_{W^{\frac{1}{2},2}\left(
\Omega\right)  }^{2}+\parallel%
%TCIMACRO{\dint \limits_{\Omega}}%
%BeginExpansion
{\displaystyle\int\limits_{\Omega}}
%EndExpansion
\Gamma_{\frac{1}{2}}\left(  y-x\right)  f(y)d\Omega_{y}\parallel_{W^{\frac
{1}{2},2}\left(  \Omega\right)  }^{2}\\
& =\parallel\left[  u\right]  _{g}\parallel_{W^{\frac{1}{2},2}\left(
\Omega\right)  }^{2}+\parallel\left[  u\right]  _{f}\parallel_{W^{\frac{1}%
{2},2}\left(  \Omega\right)  }^{2}%
\end{align*}


$\left(  iii\right)  $ \ The third result follows from the fact that
\begin{align*}
\left\langle u,\text{ }\left[  u\right]  _{f}\right\rangle _{W^{\frac{1}{2}%
,2}\left(  \Omega\right)  }  & =\left\langle
%TCIMACRO{\dint \limits_{\partial\Omega}}%
%BeginExpansion
{\displaystyle\int\limits_{\partial\Omega}}
%EndExpansion
\Gamma_{\frac{1}{2}}\left(  y-x\right)  \upsilon_{\frac{1}{2}}\left(
y\right)  g(y)d\partial\Omega_{y}+\left(  -%
%TCIMACRO{\dint \limits_{\Omega}}%
%BeginExpansion
{\displaystyle\int\limits_{\Omega}}
%EndExpansion
\Gamma_{\frac{1}{2}}\left(  y-x\right)  f(y)d\Omega_{y}\right)  ,-%
%TCIMACRO{\dint \limits_{\Omega}}%
%BeginExpansion
{\displaystyle\int\limits_{\Omega}}
%EndExpansion
\Gamma_{\frac{1}{2}}\left(  y-x\right)  f(y)d\Omega_{y}\right\rangle
_{W^{\frac{1}{2},2}\left(  \Omega\right)  }\\
& =\left\langle -%
%TCIMACRO{\dint \limits_{\Omega}}%
%BeginExpansion
{\displaystyle\int\limits_{\Omega}}
%EndExpansion
\Gamma_{\frac{1}{2}}\left(  y-x\right)  f(y)d\Omega_{y},\text{ }-%
%TCIMACRO{\dint \limits_{\Omega}}%
%BeginExpansion
{\displaystyle\int\limits_{\Omega}}
%EndExpansion
\Gamma_{\frac{1}{2}}\left(  y-x\right)  f(y)d\Omega_{y}\right\rangle
_{W^{\frac{1}{2},2}\left(  \Omega\right)  }\\
& =\parallel-%
%TCIMACRO{\dint \limits_{\Omega}}%
%BeginExpansion
{\displaystyle\int\limits_{\Omega}}
%EndExpansion
\Gamma_{\frac{1}{2}}\left(  y-x\right)  f(y)d\Omega_{y}\parallel_{W^{\frac
{1}{2},2}\left(  \Omega\right)  }^{2}%
\end{align*}
and likewise $\left(  iv\right)  $ follows from
\begin{align*}
\left\langle u,\text{ }\left[  u\right]  _{g}\right\rangle _{W^{\frac{1}{2}%
,2}\left(  \Omega\right)  }  & =\text{ }\left\langle
%TCIMACRO{\dint \limits_{\partial\Omega}}%
%BeginExpansion
{\displaystyle\int\limits_{\partial\Omega}}
%EndExpansion
\Gamma_{\frac{1}{2}}\left(  y-x\right)  \upsilon_{\frac{1}{2}}\left(
y\right)  g(y)d\partial\Omega_{y},\text{ }%
%TCIMACRO{\dint \limits_{\partial\Omega}}%
%BeginExpansion
{\displaystyle\int\limits_{\partial\Omega}}
%EndExpansion
\Gamma_{\frac{1}{2}}\left(  y-x\right)  \upsilon_{\frac{1}{2}}\left(
y\right)  g(y)d\partial\Omega_{y}\right\rangle _{W^{\frac{1}{2},2}\left(
\Omega\right)  }\\
& =\text{ }\parallel%
%TCIMACRO{\dint \limits_{\partial\Omega}}%
%BeginExpansion
{\displaystyle\int\limits_{\partial\Omega}}
%EndExpansion
\Gamma_{\frac{1}{2}}\left(  y-x\right)  \upsilon_{\frac{1}{2}}\left(
y\right)  g(y)d\partial\Omega_{y}\parallel_{W^{\frac{1}{2},2}\left(
\Omega\right)  }^{2}%
\end{align*}
.
\end{proof}


\end{document}
